\subsection{Localized acceptance}\label{sec:localized}
The acceptance rule of Proposition~\ref{prop:accept} needs a certified gap
below $1/(\Omega W)$, which on small instances already means gaps of
$2^{-50}$ to $2^{-350}$ (Section~\ref{sec:computation}). The filtering history
permits a different test that accepts a candidate as soon as the retained box
is small compared with a first-order margin. In this subsection $F$ is a
quadratic with gradient $\ell(x)=\nabla F(x)$, the grids use corrections
$d_i(v)=L_iw_i(v)^2/8$, and coordinates with $L_i=0$ use the two endpoints of
their current interval. Consider one stage with current box $B^\circ$,
min-marginals $m_i$, incumbent value $U$ and filtered box $B$, reached from $X$
by filtering steps with values of feasible points; by
Proposition~\ref{prop:filter}, every $x\in X$ with $F(x)\le U$ lies in $B$.

\begin{lemma}[Node exclusion]\label{lem:node}
For every coordinate $i$, node $v\in G_i$ and $x\in B^\circ\cap X$ with $x_i=v$,
$F(x)\ge m_i(v)$. Consequently, for $x\in B^\circ\cap X$ with $F(x)\le U$:
\begin{enumerate}[label=(\alph*)]
\item if $i\in I_Z$, then $x_i$ lies in the set $K_i$ of nodes $v$ with
$m_i(v)\le U$ and of integers strictly inside grid intervals $[a,a']$ with
$a'-a\ge2$ and $\min\{m_i(a),m_i(a')\}\le U$;
\item if $L_i=0$, $G_i=\{a,a'\}$ and $m_i(a)>U$, then setting $x_i=a'$ gives a
point of $B^\circ\cap X$ with value at most $F(x)$.
\end{enumerate}
\end{lemma}

\begin{proof}
In the proof of Proposition~\ref{prop:cellwise}(a), coordinate $i$ of $x$ is
already a vertex coordinate of every cell with $C_i\ni v$, so it need not be
moved; the resulting vertex $y$ has $y_i=v$ and
$F(x)\ge F(y)-\sum_kL_k\Delta_k(C_k)^2/8\ge Q(y)\ge m_i(v)$. Part~(a) follows from
this and Proposition~\ref{prop:cellwise}(c). For (b), $F$ is concave along
coordinate $i$, so $F(x)$ is at least the smaller of its values at $x_i=a$ and
$x_i=a'$, and the first is at least $m_i(a)>U\ge F(x)$.
\end{proof}

\begin{proposition}[Localized exact acceptance]\label{prop:local}
Define a box $B'\subseteq B$ coordinatewise: for $i\in I_Z$ with $L_i>0$ let
$B'_i$ be the hull of $K_i\cap B_i$; for $i$ with $L_i=0$ and grid $\{a,a'\}$ let
$B'_i=\{a'\}$ if $m_i(a)>U\ge m_i(a')$, $B'_i=\{a\}$ if $m_i(a')>U\ge m_i(a)$, and
$B'_i=B_i$ otherwise; and let $B'_i=B_i$ for the remaining coordinates. Let
$\hat x\in B'\cap X$ with $F(\hat x)\le U$, $\zeta=\nabla F(\hat x)$,
$W=\{i:B'_i\text{ has positive length}\}$ and $B'_i=[\alpha_i,\alpha_i']$ for
$i\in W$. Suppose every $i\in W$ falls under one of the cases
\[
\begin{array}{lll}
\text{(i)}&\hat x_i=\alpha_i,\ \zeta_i\ge0:&\mu_i=\zeta_i/(\alpha_i'-\alpha_i);\\
\text{(ii)}&\hat x_i=\alpha_i',\ \zeta_i\le0:&\mu_i=-\zeta_i/(\alpha_i'-\alpha_i);\\
\text{(iii)}&\alpha_i<\hat x_i<\alpha_i',\ \zeta_i=0:&\mu_i=0;\\
\text{(iv)}&i\in I_Z\text{ otherwise}:&\mu_i=-\abs{\zeta_i},
\end{array}
\]
and that $H_{WW}+2\diag(\mu_W)\succeq0$. Then $\hat x$ is a global minimizer.
\end{proposition}

\begin{proof}
Let $x\in X$. If $F(x)>U$, then $F(x)>F(\hat x)$. Otherwise $x\in B$. Applying
Lemma~\ref{lem:node}(b) to the coordinates with $L_i=0$ whose $B'_i$ is a
single endpoint keeps the point in $B\cap X$, does not increase $F$, and gives
a point $x''$ that, by Lemma~\ref{lem:node}(a), lies in $B'\cap X$. Let
$d=x''-\hat x$; then $d_i=0$ for $i\notin W$ and
$\alpha_i-\hat x_i\le d_i\le\alpha_i'-\hat x_i$ for $i\in W$. In cases (i) and (ii),
$\zeta_id_i=\abs{\zeta_i}\abs{d_i}\ge\mu_id_i^2$; in case (iii) both sides vanish;
in case (iv) $d_i\in\Z$, so $\zeta_id_i\ge-\abs{\zeta_i}\abs{d_i}\ge-\abs{\zeta_i}d_i^2$.
Hence
$F(x'')-F(\hat x)=\zeta^\T d+\frac12d^\T Hd\ge\frac12d_W^\T\bigl(H_{WW}+2\diag(\mu_W)\bigr)d_W\ge0$,
and $F(x)\ge F(x'')\ge F(\hat x)$.
\end{proof}

The test needs one exact $LDL^\T$ factorization of a $\abs W\times\abs W$ matrix,
which a checker repeats after replaying the filtering history; the candidate
is arbitrary, for example the stationary point of the face of a feasible
incumbent. Validity needs neither uniqueness nor growth. Under growth and
strict complementarity the test succeeds after a number of stages governed by
a geometric margin rather than by the height bound.

\begin{corollary}[Acceptance after finitely many stages]\label{cor:local}
Let $x^*$ be the unique minimizer with growth constant $g$, let $C=\{i:L_i=0\}$,
$Z=I_Z\setminus C$, $S=\{i\in I_C:\ell_i<x^*_i<u_i\}$ and
$A=\{i\in I_C\setminus C:x^*_i\in\{\ell_i,u_i\}\}$, and assume strict
complementarity, $\partial_iF(x^*)\ne0$ for $i\in A$. Put
$\lambda_A=\min_{i\in A}\abs{\partial_iF(x^*)}$,
$\Gamma=\norm{H_{AA}}_2+\norm{H_{SA}}_2^2/(2g)$, $\delta_C=\min_{i\in C}(u_i-\ell_i)$,
and
\[
h^*=\frac1{\sqrt{n\kappa}}\min\Bigl\{\tfrac12\,[Z\ne\emptyset],\ \
\tfrac45\,\delta_C\,[C\ne\emptyset],\ \ \frac{\lambda_A}{5\Gamma}\,[A\ne\emptyset]\Bigr\},
\]
where a bracketed term is omitted when its condition fails. In a trial of the
common-mesh variant with $8\kappa\theta^2\le1$, at every stage $j$ with
$h_j\le h^*$ the test of Proposition~\ref{prop:local} accepts $\hat x=x^*$. Since
$\lambda_A\ge1/(\Delta R)$ and $\delta_C\ge1/\Delta$ for rational data, this
happens after $\poly(I)+O(\log\kappa)$ stages.
\end{corollary}

The proof is in Appendix~\ref{app:localized}. In the experiments of
Section~\ref{sec:computation} the test accepted 29 of 30 random unplanted
mixed-integer instances within four stages, whereas the height rule needed up
to 542 stages. The idea of certifying a neighbourhood of a candidate by local
information is that of exclusion boxes in interval branch-and-bound
\cite{Neumaier2004}; here it is carried out exactly on the box produced by the
filtering history. With several minimizers the test can fail, for example when
two optimal vertices differ in a concave coordinate.
